
@@M0@@

\subsection*{3.4.}
Supposons enfin que les morphismes (3.2.3) et (3.3.1) soient tous les deux définis pour un morphisme @@M1@@ tel que @@M2@@. C'est le cas par exemple si @@M3@@ est quasi-compact avec un @@M4@@-module ample et si @@M5@@, où @@M6@@ est un @@M7@@-module localement libre de rang @@M8@@. Alors la formule de projection dans la catégorie dérivée donne, pour tout @@M9@@ et @@M10@@, la formule de projection
@@M11@@
Lorsque @@M12@@ est noethérien, alors 1.3.2 permet d'en déduire la formule graduée correspondante.

\section*{4. Nombres d'intersection}

\subsection*{4.1.}
Considérons un schéma propre @@M13@@ sur un corps @@M14@@. En composant le produit (3.1.1)
@@M15@@
avec la caractéristique d'Euler--Poincaré
@@M16@@
on obtient un accouplement appelé «~\emph{nombre d'intersection}~» :
@@M17@@
